\chapter{Ciała i rozszerzenia}

\section{Ciała, zanurzenia i charakterystyka}

\begin{definicja}
Ciało $K$ to zbiór z dodawaniem i mnożeniem, w którym $(K,+)$ oraz $(K^\times,\cdot)$ są grupami abelowymi, mnożenie jest rozdzielne względem dodawania i $0\ne1$. Tutaj $K^\times=K\setminus\{0\}$. Homomorfizm ciał zachowuje $0,1,+,\cdot$. Rozszerzenie $L/K$ oznacza, że $K$ jest podciałem $L$.
\end{definicja}

Homomorfizm ciał jest różnowartościowy: gdyby $a\ne0$ miał obraz zero, to obraz $aa^{-1}=1$ również byłby zerem. Dlatego nazywamy go \emph{zanurzeniem}. Zanurzenie $L\to M$ jest zanurzeniem \emph{nad $K$}, jeżeli ustala każdy element wspólnego podciała $K$.

\begin{przyklad}
Ciałami są $\Q,\R,\C$, ciało funkcji wymiernych $K(t)$ oraz $\F_p=\Z/p\Z$ dla liczby pierwszej $p$. Pierścień $\Z$ nie jest ciałem: liczba $2$ nie ma w nim odwrotności. Symbol $K(\alpha_1,\dots,\alpha_r)$ oznacza najmniejsze podciało zawierające $K$ i wskazane elementy.
\end{przyklad}

Charakterystyka $K$ jest najmniejszym $n>0$, dla którego $n\cdot1=0$, lub zerem, jeśli takie $n$ nie istnieje. Dodatnia charakterystyka musi być pierwsza: $n=ab$ z $1<a,b<n$ dawałoby iloczyn dwóch niezerowych elementów równy zeru. W charakterystyce $p$ wzór Newtona daje $(a+b)^p=a^p+b^p$. Odwzorowanie $a\mapsto a^p$ jest różnowartościowym endomorfizmem, zwanym endomorfizmem Frobeniusa. Jest automorfizmem na ciele skończonym; na dowolnym ciele charakterystyki $p$ nie musi być „na”.

\noindent\textbf{Konwencja.} Dowody dotyczące zanurzeń, teorii Galois i pierwiastników prowadzimy dalej w charakterystyce zero. Osobny przykład ciał skończonych będzie wyraźnie oznaczony. Dzięki temu nie trzeba wprowadzać teorii rozszerzeń nierozdzielczych.

\section{Dołączenie pierwiastka i wielomian minimalny}

Wielomian dodatniego stopnia z $K[T]$ jest \emph{nierozkładalny}, jeśli nie jest iloczynem dwóch wielomianów dodatniego stopnia z $K[T]$. Korzystamy z dzielenia z resztą i algorytmu Euklidesa w $K[T]$.

\begin{twierdzenie}[Dołączenie formalnego pierwiastka]
Jeżeli $p\in K[T]$ jest nierozkładalny i $d=\deg p$, to $K[T]/(p)$ jest ciałem. Element $\alpha=T+(p)$ spełnia $p(\alpha)=0$, a każdy element tego ciała ma jednoznaczną postać
\[
a_0+a_1\alpha+\cdots+a_{d-1}\alpha^{d-1},\qquad a_j\in K.
\]
\end{twierdzenie}
\begin{proof}
Dzielenie przez $p$ daje reprezentanta stopnia mniejszego niż $d$. Dwaj różni tacy reprezentanci nie mogą różnić się o wielokrotność $p$, co dowodzi jednoznaczności. Dla niezerowej klasy wybierzmy reprezentanta $q$ niepodzielnego przez $p$. Nierozkładalność $p$ daje $\gcd(p,q)=1$, więc istnieją $u,v\in K[T]$ z $up+vq=1$. Klasa $v$ jest odwrotnością klasy $q$. Stałe z $K$ zachowują różne klasy, więc $K$ zanurza się w ilorazie. Wreszcie klasa $p(T)$ jest zerowa, czyli $p(\alpha)=0$.
\end{proof}

Element $\alpha$ jest \emph{algebraiczny nad $K$}, jeśli zeruje niezerowy wielomian z $K[T]$; w przeciwnym razie jest \emph{przestępny}. Dla algebraicznego $\alpha$ wybieramy unormowany wielomian $m_{\alpha,K}$ najmniejszego stopnia zerujący $\alpha$. Dzielenie z resztą pokazuje, że dzieli on każdy wielomian zerujący $\alpha$. Jest nierozkładalny, bo z $uv(\alpha)=0$ wynikałoby $u(\alpha)=0$ albo $v(\alpha)=0$, wbrew minimalności stopnia. Odwzorowanie $q\mapsto q(\alpha)$ daje więc
\[
K[T]/(m_{\alpha,K})\simeq K[\alpha]=K(\alpha).
\]
Ostatnia równość wynika z tego, że iloraz jest już ciałem, a nie tylko pierścieniem. Jest to uzasadnienie wszystkich relacji między potęgami $\alpha$, którego będziemy potrzebować przy przedłużaniu zanurzeń.

\begin{przyklad}
$\Q[T]/(T^2-2)\simeq\Q(\sqrt2)$ oraz $\R[T]/(T^2+1)\simeq\C$. W pierwszym przykładzie każdy element ma postać $a+b\sqrt2$, $a,b\in\Q$.
\end{przyklad}

\section{Stopień rozszerzenia i twierdzenie o wieży}

Stopień $[L:K]$ to wymiar przestrzeni wektorowej $L$ nad $K$. Z konstrukcji ilorazowej wynika $[K(\alpha):K]=\deg m_{\alpha,K}$, z bazą $1,\alpha,\dots,\alpha^{d-1}$.

\begin{twierdzenie}[O wieży]
Jeśli $K\subseteq E\subseteq L$, a $[E:K]$ i $[L:E]$ są skończone, to
\[
[L:K]=[L:E]\,[E:K].
\]
Jeśli $[L:K]$ jest skończony, oba stopnie po prawej także są skończone i zachodzi ta sama równość.
\end{twierdzenie}
\begin{proof}
Niech $u_1,\dots,u_r$ będzie bazą $E/K$, a $v_1,\dots,v_s$ bazą $L/E$. Każde $x\in L$ zapisujemy jako $\sum_j b_jv_j$, a następnie każde $b_j$ jako $\sum_i a_{ij}u_i$ z $a_{ij}\in K$. Iloczyny $u_iv_j$ rozpinają więc $L$ nad $K$. Jeżeli $\sum_{i,j}a_{ij}u_iv_j=0$, to niezależność $v_j$ nad $E$ daje $\sum_i a_{ij}u_i=0$ dla każdego $j$. Niezależność $u_i$ nad $K$ daje wszystkie $a_{ij}=0$. Iloczyny tworzą bazę złożoną z $rs$ elementów.

Jeśli $[L:K]<\infty$, to $E$ jest podprzestrzenią skończonego wymiaru. Ponadto każda rodzina niezależna nad $E$ jest niezależna nad $K$, więc $[L:E]$ też jest skończony. Stosujemy już udowodnioną część.
\end{proof}

Każdy element skończonego rozszerzenia jest algebraiczny: dostatecznie wiele jego potęg jest liniowo zależnych. Odwrotnie, dołączenie skończenie wielu elementów algebraicznych daje rozszerzenie skończone, na mocy twierdzenia o wieży. Nie oznacza to, że każde rozszerzenie algebraiczne ma skończony stopień.

\section{Przedłużanie zanurzeń: podstawowe narzędzie}

Ciało $\Omega$ jest \emph{algebraicznie domknięte}, gdy każdy wielomian dodatniego stopnia z $\Omega[T]$ ma w $\Omega$ zero. Domknięcie algebraiczne $\overline K/K$ jest rozszerzeniem algebraicznym, którego ciało $\overline K$ jest algebraicznie domknięte. Korzystamy ze standardowego twierdzenia o istnieniu domknięcia algebraicznego, którego konstrukcja teoriomnogościowa nie jest tu potrzebna. W przykładach liczbowych możemy umieszczać wszystkie rozważane rozszerzenia w $\C$.

\begin{lemat}[Przedłużenie przez jeden element]
Niech $\sigma:K\to\Omega$ będzie zanurzeniem, a $\alpha$ algebraiczny nad $K$, z $p=m_{\alpha,K}$. Dla każdego zera $\beta$ wielomianu $\sigma(p)$, otrzymanego przez zastosowanie $\sigma$ do współczynników, istnieje dokładnie jedno przedłużenie
\[
\widetilde\sigma:K(\alpha)\longrightarrow\Omega,
\qquad \widetilde\sigma(\alpha)=\beta.
\]
Każde przedłużenie powstaje w ten sposób.
\end{lemat}
\begin{proof}
Wielomian $\sigma(p)$ jest nierozkładalny nad $\sigma(K)$: jego rozkład po zastosowaniu $\sigma^{-1}$ dałby rozkład $p$. Wobec tego jest wielomianem minimalnym $\beta$ nad $\sigma(K)$. Wzór
\[
\sum_{j=0}^{d-1}a_j\alpha^j\longmapsto
\sum_{j=0}^{d-1}\sigma(a_j)\beta^j
\]
indukuje izomorfizm odpowiednich ilorazów przez $(p)$ i $(\sigma(p))$, a następnie zanurzenie do $\Omega$. Zachowuje więc wszystkie działania i relacje, nie tylko relację $p(\alpha)=0$. Jednoznaczność wynika z generowania ciała przez $K$ i $\alpha$. Każde zanurzenie musi wysłać $\alpha$ na zero $\sigma(p)$, bo zachowuje równania wielomianowe.
\end{proof}

\begin{twierdzenie}[Przedłużanie zanurzeń]
Niech $L/K$ będzie rozszerzeniem skończonym, a $\Omega$ ciałem algebraicznie domkniętym. Każde zanurzenie $K\to\Omega$ przedłuża się do zanurzenia $L\to\Omega$.
\end{twierdzenie}
\begin{proof}
Wybierzmy skończoną rodzinę generatorów $L=K(\alpha_1,\dots,\alpha_r)$, na przykład bazę $L/K$. Po skonstruowaniu zanurzenia $K(\alpha_1,\dots,\alpha_{i-1})$ stosujemy poprzedni lemat do wielomianu minimalnego $\alpha_i$ nad tym ciałem. Jego obraz ma zero w $\Omega$, więc następny krok jest możliwy. Po $r$ krokach otrzymujemy wymagane zanurzenie.
\end{proof}

\begin{twierdzenie}[Liczba zanurzeń]
Skończone rozszerzenie $L/K$ w charakterystyce zero ma dokładnie $[L:K]$ różnych $K$-zanurzeń do ustalonego domknięcia algebraicznego $\overline K$.
\end{twierdzenie}
\begin{proof}
Nierozkładalny wielomian $p$ w charakterystyce zero ma różne zera: jego pochodna jest niezerowa i ma mniejszy stopień, więc $\gcd(p,p')=1$. W wieży z poprzedniego dowodu oznaczmy kolejne stopnie przez $d_i$. Lemat 1.6 daje dokładnie $d_i$ przedłużeń każdego dotychczasowego zanurzenia, bo obraz wielomianu minimalnego ma $d_i$ różnych zer. Różne wybory dają różne zanurzenia, a każde zanurzenie powstaje z takich wyborów. Ich liczba to $\prod_i d_i=[L:K]$.
\end{proof}

\begin{przyklad}
$\Q(\sqrt2)$ ma dwa $\Q$-zanurzenia i oba są automorfizmami tego ciała. $\Q(\sqrt[3]{2})$ ma trzy $\Q$-zanurzenia do $\C$, lecz tylko jedno jest jego automorfizmem: dwa pozostałe wysyłają rzeczywisty generator na $\omega\sqrt[3]{2}$ lub $\omega^2\sqrt[3]{2}$, gdzie $\omega=e^{2\pi i/3}$. Zanurzenie wysyła zatem element na jedno z jego sprzężeń, niekoniecznie na element różny od niego.
\end{przyklad}


